\section{A source-free angular budget for the finite matched-Bob corner}
\label{sec:new-c144-matched-budget}

This is an exact auxiliary identity, not by itself a C14.4 admission.
Work in the separated tensor ring $K_c\otimes_{\mathbb R}K_t$.
Let $S,\kappa,\nu,\ell$ be positive source quantities, with
$\nu+\ell=1$, and put
\[
 x=\kappa S,\qquad j=\sqrt{\nu\ell},\qquad
 a_B=\sqrt{\ell/\nu}.
 \tag{MB1}
\]
For independent testing variables $u,v,W,Z$ define
\[
 T=\kappa W,\quad U=uZ,\quad P=uW/\sqrt S,\quad Q=\sqrt S Z,
 \quad N=T-jU,\quad G=P-jQ,
 \quad \delta=v-a_B,\quad B=\delta Q,\quad B_u=\delta U.
 \tag{MB2}
\]
The monomial relation $TU=\kappa PQ$ is exact. It gives
\[
 \delta^2N^2+2j\kappa\delta GB+2j^2\kappa B^2-j^2B_u^2
       =(\delta T)^2.
 \tag{MB3}
\]
In particular, if $\nu,\ell$ are positive bounded units,
$\kappa$ is infinitesimal and $\delta$ has zero real residue,
then any fixed real positive diagonal on $B_u$ can be purchased
at an infinitesimal perturbation of the Gram matrix on $(N,G,B)$.
Indeed, for any fixed $c>0$ the exact SOS remainder is
\[
 \frac{c}{j^2}\delta^2N^2+
       \frac{2c\kappa}{j}\delta GB+2c\kappa B^2-cB_u^2
       =\frac{c}{j^2}(\delta T)^2.
 \tag{MB4}
\]
Every displayed coefficient is a finite sum of bounded separated
factors, with zero real residue except the constant $c$ in the
desired $B_u$ pivot. The only inverses are of positive source units.
The factor $\delta$ is a difference of a pure testing and a pure
source quantity; no inverse or radical of that difference is taken.
The right side has the displayed separated SOS root
$\sqrt c\,\delta T/j$.

Two exact radial conductors and one cross conductor are
\[
 \begin{split}
 (u^2-x)W/\sqrt S&=uG-\sqrt S N,\\
 (u^2-x)Z&=(\kappa\sqrt S/j)G-(u/j)N,\\
 (u^2-x)\delta Z^2&=(\kappa/j)GB-(1/j)NB_u.
 \end{split}
 \tag{MB5}
\]
These identities require no bound on the state coordinates or on
any ratio of source and testing infinitesimals. In the local matched
chart with $u$ infinitesimal their first two right sides have only
infinitesimal coefficients on $(N,G)$.

For comparison, the four-term contact circuit is
\[
 \begin{split}
 \Psi={}&(\nu/\ell)v^2T^2+\nu^2v^2U^2
       +\kappa\nu\ell Q^2+\kappa P^2-4\kappa\nu vPQ\\
   ={}&(v/a_B)^2N^2+\kappa G^2
           +(2\kappa\nu/a_B)PQ(v-a_B)^2.
 \end{split}
 \tag{MB6}
\]
On the positive testing/state sign chart the last summand is itself
a separated square. The budget MB3 is stronger for the present local
purpose: it uses only the actual monomial relation, and does not
divide by a weak circuit coefficient.

The full scalar and all original state normals still have to be
placed in this frame before using MB4. No scalar or whole-parent
conclusion is asserted here.
